% BT973 -- rail generation and phase bookkeeping theorem insert

\subsection{Rail generation faces and phase bookkeeping (BT962--BT973)}\label{subsec:bt973-rail-generation-phase}

The final selector supplies four hyperbolic rails.  Let their support sums and
xor masks be
\[
  s=(12,12,14,22),\qquad x=(71,46,91,234).
\]
A point of the nonzero chain shadow has a support set of active rails.  The
exactly-three-rail faces are
\[
  (0,1,2),\quad (0,1,3),\quad (0,2,3),\quad (1,2,3),
\]
and each has \(27\) elements.  The high-support rail is rail \(3\).  Therefore
the three faces containing rail \(3\),
\[
  (0,1,3),\qquad (0,2,3),\qquad (1,2,3),
\]
form a selector-fixed \(27+27+27\) decomposition, while \((0,1,2)\) is the
complementary \(27\)-face.

For phase bookkeeping define the selector phase score
\[
  P(r)=s(r)+\operatorname{popcount}(x(r)).
\]
This gives
\[
  P=(16,16,19,27).
\]
The two light rails \(0\) and \(1\) remain tied by support and phase score.  The
ABI order breaks only the packet-ordering tie by ascending xor mask, so rail
\(1\) precedes rail \(0\) because \(46<71\).  This establishes a canonical
coordinate system for downstream generation and phase tests without asserting a
numerical CKM or PMNS prediction.

\begin{theorem}[Selector-fixed rail bookkeeping]
The final E8 selector canonically fixes:
\begin{enumerate}
\item a \(27+27+27\) rail-face decomposition \((0,1,3),(0,2,3),(1,2,3)\);
\item a complementary \(27\)-face \((0,1,2)\);
\item phase-slot scores \((16,16,19,27)\);
\item an ABI light-rail order \(1<0\) induced by xor masks \(46<71\).
\end{enumerate}
These are gauge-fixing data for downstream representation, generation, and
mixing tests.  They are not by themselves fitted Standard Model parameters.
\end{theorem}
